\subsection{\texorpdfstring{同态 \(\operatorname{Ext}_{\mathbf{A}}(\mathbf{P}, \mathbf{Q}) \otimes \operatorname{Tor}^{\mathbf{A}}(\mathbf{P}, \mathbf{M}) \to \operatorname{Tor}^{\mathbf{A}}(\mathbf{Q}, \mathbf{M})\)}{同态 \textbackslash operatorname\{Ext\}\_\{\textbackslash mathbf\{A\}\}(\textbackslash mathbf\{P\}, \textbackslash mathbf\{Q\}) \textbackslash otimes \textbackslash operatorname\{Tor\}\^{}\{\textbackslash mathbf\{A\}\}(\textbackslash mathbf\{P\}, \textbackslash mathbf\{M\}) \textbackslash 到 \textbackslash operatorname\{Tor\}\^{}\{\textbackslash mathbf\{A\}\}(\textbackslash mathbf\{Q\}, \textbackslash mathbf\{M\})}}

设 M 为左 A-模，P 与 Q 为两个右 A-模。考虑同态 \(\operatorname{Homgr}_{\mathbf{A}}(\mathrm{L}(\mathrm{P}), \mathrm{L}(\mathrm{Q})) \otimes_{k} (\mathrm{L}(\mathrm{P}) \otimes_{\mathbf{A}} \mathrm{L}(\mathrm{M})) \to \mathrm{L}(\mathrm{Q}) \otimes_{\mathbf{A}} \mathrm{L}(\mathrm{M})\) 它将 \(f \otimes (x \otimes y)\) 映至 \(f(x) \otimes y\) 。根据X，第99页，这是复形的一个态射。由此可推出一个 0 次的分次 \(k\) -线性映射

\[
\mathrm{H} (\operatorname{Homgr} _ {\mathrm{A}} (\mathrm{L} (\mathrm{P}), \mathrm{L} (\mathrm{Q}))) \otimes_ {k} \operatorname{Tor} ^ {\mathrm{A}} (\mathrm{P}, \mathrm{M}) \rightarrow \operatorname{Tor} ^ {\mathrm{A}} (\mathrm{Q}, \mathrm{M})
\]

因此，经由 § 6（X，第100页，定理1）的同构 \(\varphi(\mathbf{L}(\mathbf{P}),\mathbf{L}(\mathbf{Q}))\) ，得到一个 0 次的分次 \(k\) -线性映射

\[
\operatorname{Ext} _ {\mathbf {A}} (\mathrm{P}, \mathrm{Q}) \otimes_ {k} \operatorname{Tor} ^ {\mathbf {A}} (\mathrm{P}, \mathrm{M}) \rightarrow \operatorname{Tor} ^ {\mathbf {A}} (\mathrm{Q}, \mathrm{M}),\tag{10}
\]

对应于k-双线性映射

\[
c _ {\mathrm{P}, \mathrm{Q}; \mathrm{M}}: \operatorname{Ext} _ {\mathrm{A}} ^ {n} (\mathrm{P}, \mathrm{Q}) \times \operatorname{Tor} _ {m} ^ {\mathrm{A}} (\mathrm{P}, \mathrm{M}) \rightarrow \operatorname{Tor} _ {m - n} ^ {\mathrm{A}} (\mathrm{Q}, \mathrm{M});\tag{11}
\]

偶 \((\alpha, \gamma)\) 在 \(c_{P,Q;M}\) 下的像称为 \(\alpha\) 与 \(\gamma\) 的复合积，记作 \(\alpha \circ \gamma\) .

按构造， \(\alpha \circ \gamma\) 如下获得： \(\alpha\) 由复形的态射 \(f: \mathrm{L}(\mathrm{P}) \to \mathrm{L}(\mathrm{Q})(-n)\) 表示， \(\gamma\) 由元素 \(z \in Z_m(\mathrm{L}(\mathrm{P}) \otimes_{\mathbf{A}} \mathrm{L}(\mathrm{M}))\) 表示，则 \(\alpha \otimes \gamma\) 是元素

\[
(f \otimes 1) (z) \in Z _ {m} (\mathrm{L} (\mathrm{Q}) (- n) \otimes_ {\mathrm{A}} \mathrm{L} (\mathrm{M})) = Z _ {m - n} (\mathrm{L} (\mathrm{Q}) \otimes_ {\mathrm{A}} \mathrm{L} (\mathrm{M})).
\]

例如，如果 \(\alpha \in \operatorname{Hom}_{\mathbf{A}}(\mathbf{P},\mathbf{Q})\) ，则 \(\alpha \circ \gamma = \operatorname{Tor}(\alpha,1)(\gamma)\) .

注记. --- 1) 若使用 X，第 69 页的同构 \(\psi\) ，则也可通过交换图定义复合积

\[
\begin{array}{c} \operatorname{Ext} _ {\mathrm{A}} ^ {n} (\mathrm{P}, \mathrm{Q}) \times \operatorname{Tor} _ {m} ^ {\mathrm{A}} (\mathrm{P}, \mathrm{M}) \xrightarrow {c _ {\mathrm{P} , \mathrm{Q} : \mathrm{M}}} \operatorname{Tor} _ {m - n} ^ {\mathrm{A}} (\mathrm{Q}, \mathrm{M}) \\ \bar {a} _ {\mathrm{P}, \mathrm{Q}} \times \psi_ {\mathrm{P}} (\mathrm{M}) \Biggl \downarrow \\ \mathrm{H} ^ {n} (\operatorname{Homgr} _ {\mathrm{A}} (\mathrm{L} (\mathrm{P}), \mathrm{L} (\mathrm{Q}))) \times \mathrm{H} _ {m} (\mathrm{L} (\mathrm{P}) \otimes_ {\mathrm{A}} \mathrm{M}) \longrightarrow \mathrm{H} _ {m - n} (\mathrm{L} (\mathrm{Q}) \otimes_ {\mathrm{A}} \mathrm{M}); \end{array}
\]

换言之，我们将 \(\alpha\) 表示为态射 \(f\) ：从 L(P) 到 L(Q) \((-n)\) ，将 \(\gamma\) 表示为循环 \(x \in \mathrm{L}_{m}(\mathrm{P}) \otimes_{\mathbf{A}} \mathrm{M}\) ，且 \(\alpha \circ \gamma\) 是下述元素的循环类

\[
\left(f _ {m} \otimes 1 _ {\mathbf {M}}\right) (x) \in \mathrm{L} _ {m - n} (\mathrm{Q}) \otimes_ {\mathbf {A}} \mathrm{M}.
\]

\begin{enumerate}
\def\labelenumi{\arabic{enumi})}
\setcounter{enumi}{1}
\tightlist
\item
  也可以使用分解 I(P) 和 I(Q)。
\end{enumerate}

类似地，如果 N 是第二个左 A-模，则可定义复合积 \((\mu, \gamma) \mapsto \mu \circ \gamma\) ，记作

\[
c _ {\mathrm{P}; \mathrm{M}, \mathrm{N}}: \operatorname{Ext} _ {\mathrm{A}} ^ {r} (\mathrm{M}, \mathrm{N}) \times \operatorname{Tor} _ {m} ^ {\mathrm{A}} (\mathrm{P}, \mathrm{M}) \rightarrow \operatorname{Tor} _ {m - r} ^ {\mathrm{A}} (\mathrm{P}, \mathrm{N})
\]

由交换图

\[
\begin{array}{c} \operatorname{Ext} _ {\mathbf {A}} ^ {r} (\mathbf {M}, \mathbf {N}) \times \operatorname{Tor} _ {m} ^ {\mathbf {A}} (\mathbf {P}, \mathbf {M}) \xrightarrow {c _ {\mathrm{P} ; \mathrm{M} , \mathrm{N}}} \operatorname{Tor} _ {m - r} ^ {\mathbf {A}} (\mathbf {P}, \mathbf {N}) \\ 1 \times \sigma_ {\mathrm{P}, \mathrm{M}, r} \Biggl \downarrow \\ \operatorname{Ext} _ {\mathbf {A} ^ {\circ}} ^ {r} (\mathbf {M} ^ {\circ}, \mathbf {N} ^ {\circ}) \times \operatorname{Tor} _ {m} ^ {\mathbf {A} ^ {\circ}} (\mathbf {M} ^ {\circ}, \mathbf {P} ^ {\circ}) \xrightarrow {c _ {\mathrm{M} ^ {\circ} , \mathrm{N} ^ {\circ} ; \mathrm{P} ^ {\circ}}} \operatorname{Tor} _ {m - r} ^ {\mathbf {A} ^ {\circ}} (\mathbf {N} ^ {\circ}, \mathbf {P} ^ {\circ}) \end{array}\tag{12}
\]

其中 σ 表示交换同构（X，第 71 页）。

若 \(\mu \in \operatorname{Ext}_{\mathbf{A}}^{r}(\mathbf{M}, \mathbf{N})\) 是态射 \(g: \mathrm{L}(\mathbf{M}) \to \mathrm{L}(\mathbf{N})(-r)\) 的类，且 \(\gamma \in \operatorname{Tor}_{m}^{\mathbf{A}}(\mathbf{P}, \mathbf{M})\) 是 \(z = \sum_{i,j} z_{ij}\) 的循环类，其中 \(z_{ij} \in \mathrm{L}_i(\mathbf{P}) \otimes_{\mathbf{A}} \mathrm{L}_j(\mathbf{M})\) ，则 \(\mu \circ \gamma\) 就是循环 \(\sum_{i,j} (-1)^{ir}(1 \otimes g)(z_{ij})\) 的类 .

我们还可以将 \(\gamma\) 表示为一个循环 \(y \in \mathrm{P} \otimes \mathrm{L}_{m}(\mathbf{M})\) ，且 \(\mu \circ \gamma\) 是 \((1 \otimes g)(y) \in \mathrm{P} \otimes \mathrm{L}_{m-r}(\mathbf{M})\) 的循环类 .

命题 6. --- 设 K、M、N 为左 A-模，P、Q、R 为右 A-模，

\[
\alpha \in \operatorname{Ext} _ {\mathrm{A}} ^ {n} (\mathrm{P}, \mathrm{Q}), \beta \in \operatorname{Ext} _ {\mathrm{A}} ^ {p} (\mathrm{Q}, \mathrm{R}), \lambda \in \operatorname{Ext} _ {\mathrm{A}} ^ {r} (\mathrm{K}, \mathrm{M}), \mu \in \operatorname{Ext} _ {\mathrm{A}} ^ {s} (\mathrm{M}, \mathrm{N}), \gamma \in \operatorname{Tor} _ {m} ^ {\mathrm{A}} (\mathrm{P}, \mathrm{K}).
\]

则

\begin{enumerate}
\def\labelenumi{(\arabic{enumi})}
\setcounter{enumi}{12}
\tightlist
\item
\end{enumerate}

\[
(\beta \circ \alpha) \circ \gamma = \beta \circ (\alpha \circ \gamma) \text { in }\operatorname{Tor} _ {m - p - n} ^ {\mathrm{A}} (\mathrm{R}, \mathrm{K}),\tag{14}
\]

\[
(\mu \circ \lambda) \circ \gamma = \mu \circ (\lambda \circ \gamma) \text { in }\operatorname{Tor} _ {n - r - s} ^ {\mathbf {A}} (\mathrm{P}, \mathrm{N}),\tag{15}
\]

\[
\alpha \circ (\lambda \circ \gamma) = (- 1) ^ {n r} \lambda \circ (\alpha \circ \gamma) \text { in }\operatorname{Tor} _ {m - p - r} ^ {\mathrm{A}} (\mathrm{Q}, \mathrm{M}).
\]

公式（13）和（14）直接由定义得出。我们来证明（15）。设 \(z = \sum z_{ij}, z_{ij} \in \mathbf{L}_i(\mathbf{P}) \otimes \mathbf{L}_j(\mathbf{K})\) 为表示 \(\gamma, f: \mathrm{L}(\mathrm{P}) \to \mathrm{L}(\mathrm{Q})(-n)\) 与 \(g: \mathbf{L}(\mathbf{K}) \to \mathbf{L}(\mathbf{M})(-r)\) 的循环与态射，它们分别表示 \(\alpha\) 与 \(\lambda\) 。则 \(\lambda \circ (\alpha \circ \gamma)\) 是 \(\sum (-1)^{(i-n)r}(f \otimes g)(z_{ij})\) 的类，而 \(\alpha \circ (\lambda \circ \gamma)\) 是

\[
\sum (- 1) ^ {i r} (f \otimes g) (z _ {i j}), \quad \text { whence } (1 5).
\]

